% \section{本研究の限界} \label{sec6}
\section{おわりに} \label{sec6}
本研究では，質問に応じてRAGの構成を適応的に選択し，精度とコストのトレードオフを調整する手法を提案した．
具体的には，複数のRAG構成に対してRM学習を適用し，質問特徴，文書特徴，構成情報から有用なRAG構成を上位に提示するモデルを構築した．
その上で，適合率，再現率，$F_1$値，実行時間を重み付きで統合した評価指標に基づき，最終的な構成を選択する手法を設計した．
運用上の要求に応じて重みを変更することで，精度重視やコスト重視といった方針の切り替えが可能となった．
評価項目(ii)における性能評価では，全クエリに同一のRAG構成を適用する固定構成（標準設定・最頻出設定）や，既存の代表的な適応的RAG手法であるAdaptive-RAGと比較して，回答精度（BERTScore $F_1$値，EM，Token-level $F_1$値）の大幅な改善と実行時間の短縮（Pareto改善）を確認した．
さらに，Wilcoxon符号付き順位検定により，提案手法のベースラインに対する統計的有意な優位性を立証した．
次に，提示された上位候補からユーザが最終選択を行う評価項目(iii)では，すべての候補を試行する累積コストが発生するものの，精度がさらに向上し，時間重視方針においてはオラクルの時間最良設定と完全に一致するコスト最小化を達成した．
アブレーション分析により特徴量の重要度および提案手法全体の有効性を確認し，特徴量貢献度分析により生成・検索モデルなどのマクロ要素に加えてコサイン類似度やトークン数といったデータ依存の特徴量がランキング予測に強く寄与していることを明らかにした．

今後の課題は，RM（ランキングモデル）の予測精度のさらなる向上，および現在は手動で行っている優先度重みの自動推定メカニズムの構築である．具体的には，入力クエリの難易度やシステムの実行時制約から，重みベクトルを自動決定するゲートネットワークの設計などが挙げられる．

% 今後の課題は，RMの精度向上，ユーザ意図から重み設定の自動化，および評価指標の拡張である．

% \section{今後の課題} \label{sec6.2}
% 本研究の枠組みは有効性を示した一方で，オラクルの適合率最良構成，再現率最良構成，$F_1$値最良構成との差が残り，実用化に向けて複数の課題がある．

% 第一に，ランキングモデルの精度の向上が必要である．
% 精度指標を目的変数とした場合，生成の揺らぎや取得文書の質のばらつきにより順位が不安定になりやすい．
% そのため，学習データの拡充やモデル設計の改善に加え，複数回実行による期待値推定や不確実性を考慮した学習といった，評価のばらつきを前提とした頑健な順位学習が求められる．

% 第二に，候補設定集合の設計と探索範囲の最適化が課題である．
% 候補集合が有限である以上，理想設定が候補外に存在すれば上界に到達できない．
% その一方で，候補数が増加すると学習難度と運用負荷が増大する．
% したがって，多様性を確保しつつ冗長性を抑えた候補設計や，段階的に候補を広げる探索戦略を導入する必要がある．
% また，将来的には検索戦略，再ランキング，生成手順といった，より高次の設計要素を含めた適応化も検討課題となる．

% 第三に，重み設定の自動化とユーザ意図への適応が挙げられる．
% 本研究では重みを手動で設定した．
% しかし，実運用ではタスクや状況により要求が変化する．
% そのため，質問特性や環境制約から重みを推定する仕組み，あるいは少数のプリセットに整理したインタフェース設計が必要である．

% 第四に，上位$k$候補を用いる運用では，最終的に採用された設定の時間だけでなく，候補を複数回試行する総コストが重要となる．
% 今後は総コストを評価に含め，閾値に基づく早期終了といった段階的運用を含めた最適化が求められる．

% 第五に，評価指標の拡張が必要である．
% 本研究は適合率・再現率・$F_1$値と時間を中心に評価した．
% しかし，RAGの実運用では事実性や根拠の正確性といった，信頼性に直結する品質要因が重要である．
% 人手評価や根拠整合性評価を取り込み，目的変数として扱うことで，より実問題に即した適応的選択が可能になると考えられる．

% 以上の課題に取り組むことで，本研究の提案手法は，高精度側の上界に近づきつつ，総コスト制約下でも運用可能な適応的RAG選択手法として発展すると期待される．